% docs/physics/latex/09_phase_space.tex - SIMC++ physics note 09: phase-space (acceptance) runs, doing_phsp.
% Standalone PDF docs/physics/pdf/SIMCpp_Physics_09_phase_space.pdf (no folder copy: the reaction lives in
% src/reactions/HydrogenElastic.cpp, see notes.json); a chapter of the book SIMC++ Physics Reference.
% Rules: docs/physics/README.md; structure: template_reaction.tex. Every display cites C++ first, Fortran second.
\documentclass{simcnote}
\simccommon
\begin{document}
\simcnotemeta{09}{Phase space}{Phase-space (acceptance) runs: \texttt{doing\_phsp}}%
  {Reaction plug-in \texttt{phase\_space}: the H(e,e\pr p) generation with unit cross section and unit normalisation;
   what the switch does to every other reaction}
\simcnotemaketitle
\simcabstract{The input \code{doing\_phsp} = 1 asks SIMC for an acceptance run: the cross section is replaced by 1 and
the normalisation factor by 1. It is not a separate generator. It removes most reaction flags but keeps the target-class
flags and the hadron mass that they selected, so the events are generated by whatever reaction code those leftover
flags reach. In practice, and in SIMC++, the only working combination is hydrogen with no other reaction flag: the
events are H(e,e\pr p) elastic events (note 01) with weight $J\,w_\mathrm{gen}$. This note gives the formulas of that
run, what the switch does to each reaction in the Fortran, and what SIMC++ accepts. Notation, glossary and common
physics: note 00.}
\simcnotetoc

% =====================================================================================================================
\section{The reaction in SIMC++}\lbl{sec:reaction}

\subsection{Process and scope}
An acceptance (``phase-space'') run counts how many generated events reach the detectors, without a physics weight. In
SIMC++ \code{doing\_phsp} = 1 on hydrogen, with no other reaction flag, runs the generation and the kinematics of
H(e,e\pr p) elastic scattering: the electron direction is generated flat in its slopes, the electron energy and the proton
follow from the elastic relation, and the event weight is the solid-angle jacobian times the radiative weight. What is
simulated: everything of note 01 except the cross section, i.e.\ beam, energy loss, multiple scattering, radiation (it
stays on, \secr{sec:rad}), both spectrometers, the reconstruction. What is not: a flat phase space in the hadron
variables (the proton is fixed by the electron), any other reaction. The header of the \texttt{.hist} says
``PHASE SPACE - NO physics, NO radiation (may not work)'', and the start-up message ``phsp option is not really
implemented - buyer beware!!''; the second is accurate, the first is not for radiation.

\begin{simcwarning}
\textbf{Scope.} SIMC++ runs \code{doing\_phsp} only on hydrogen without another reaction flag. Every other combination
is refused or gives no event (\tabr{tab:replaces}): deuterium and heavier targets are a configuration error (exit
status 2), a SIDIS flag is refused (exit status 3), and with a pion, kaon, Delta or $\rho$ flag the run stops with ``no
layout'' when a tree is written (exit status 3) and no try reaches the spectrometers without a tree.
\end{simcwarning}

\subsection{Components}
\begin{simctable}{footnotesize}{L{0.14\linewidth}L{0.16\linewidth}L{0.15\linewidth}L{0.25\linewidth}L{0.30\linewidth}}
\caption{The SIMC++ components of a phase-space run (one component per row).}\lbl{tab:components}\\
\toprule\rowcolor{simchead}\simcth{role} & \simcth{interface} & \simcth{registered name} & \simcth{class} & \simcth{file}\\\midrule\endfirsthead
\toprule\rowcolor{simchead}\simcth{role} & \simcth{interface} & \simcth{registered name} & \simcth{class} & \simcth{file}\\\midrule\endhead
\bottomrule\endlastfoot
reaction & \code{IReaction} & \code{phase\_space} & \code{PhaseSpace} (a \code{HydrogenElastic}) & \file{reactions/HydrogenElastic.cpp}\\
event completion and reconstruction & (base class) & -- & \code{EepReactionBase} & \file{reactions/EepReactionBase.cpp}\\
weight assembly & (function) & -- & \code{reactions::assembleWeight} & \file{reactions/EventChain.cpp}\\
normalisation & (function) & -- & \code{Normalise} & \file{driver/RunReport.cpp}\\
reaction flags & (function) & -- & \code{DeriveReactionFlags} & \file{core/RunSetup.cpp}\\
masses, radiation switch & (function) & -- & \code{BuildReaction\-AndKinematics} & \file{setup/SetupReaction.cpp}\\
cross section & \code{IPhysicsModel} & \code{sigep} & \code{proton::SigEPModel} & \file{physics/proton/ProtonModels.cpp} (created, never called)\\
beam, target, radiation, arms, outputs & as in note 00 & \code{standard}, \code{hallc\_cryo}, \code{simc\_rad}, \ldots & & \ctabr{tab:plugins}\\
\end{simctable}

\subsection{Selection and override}
\code{doing\_phsp} = 1 makes the pipeline choose \code{phase\_space} before any other rule (\ctabr{tab:defaults}). The
flags are derived first as without the switch (the target class, the hadron mass $M_h$), then the switch clears
\code{doing\_eep}, \code{doing\_pion}, \code{doing\_kaon}, \code{doing\_delta}, \code{doing\_rho} and sets the radiation
input \code{rad\_flag} to 0 \fsrc{dbase.f:231-238}. It does not clear \code{doing\_semi}, \code{doing\_hyd\_elast},
\code{doing\_deuterium}, \code{doing\_heavy} nor the target-class flags of the meson reactions, and it does not switch
\code{using\_rad} off. Because \code{doing\_eep} is cleared before the struck-nucleon masses are set, no branch sets them:
$M_N=0$ and the recoil mass is 0 \fsrc{dbase.f:297-317}. The plug-in \code{PhaseSpace} refuses, in its \code{init()},
deuterium and heavier targets and a SIDIS flag; a meson flag has already been cleared by the switch, so such a set-up
passes \code{init()} and fails later (\tabr{tab:replaces}). A JSON input can name \code{"reaction": "phase\_space"}
directly; the conditions are the same.

\subsection{What the switch does to each reaction}
\begin{simctable}{footnotesize}{L{0.17\linewidth}L{0.5\linewidth}L{0.33\linewidth}}
\caption{\code{doing\_phsp} = 1 combined with each reaction (one combination per row). Fortran: from the code (only the
hydrogen row has Fortran reference runs). SIMC++: probe runs of this note (\secr{sec:validation}).}\lbl{tab:replaces}\\
\toprule\rowcolor{simchead}\simcth{combination} & \simcth{Fortran (59ccdf6)} & \simcth{SIMC++}\\\midrule\endfirsthead
\toprule\rowcolor{simchead}\simcth{combination} & \simcth{Fortran (59ccdf6)} & \simcth{SIMC++}\\\midrule\endhead
\bottomrule\endlastfoot
H, no other flag & elastic generation and kinematics (note 01), $\sigma=1$, normfac $=1$ & runs (this note)\\
D, no other flag & D(e,e\pr p) generation; with $M_N=0$ the residual mass becomes $M_p+M_n$, the proton energy is solved for that recoil; $S=1$, $\sigma=1$ & refused at set-up (configuration error, exit status 2)\\
$A\ge3$, no other flag & A(e,e\pr p) generation; with $M_N=0$ the vertex $E_m=M^*_\mathrm{rec}-M_A$ is negative and fails the vertex limit $E_m\ge E_\mathrm{Fermi}$: no event & refused at set-up (exit status 2)\\
\code{doing\_nuc\_elast} & the elastic flag is kept & refused at set-up\\
\code{doing\_pion} & flag cleared, target class kept: only the electron slopes are generated, no energy; the hadron is put at $P_h^0$ along its previous slopes & runs only without tree (\code{Nntu} = 0): no try reaches the spectrometers, \texttt{wtcontr} $=0$; with a tree: stop ``ntuple: no layout for this reaction'' (exit status 3)\\
\code{doing\_kaon} & as \code{doing\_pion} & as \code{doing\_pion}\\
\code{doing\_delta} & as \code{doing\_pion} & as \code{doing\_pion}\\
\code{doing\_rho} & as \code{doing\_pion} (no $\rho$ decay) & as \code{doing\_pion}\\
\code{doing\_semi}, pions & SIDIS generation and kinematics ($M_N=M_p$), $\sigma=1$, normfac $=1$, SIDIS tree layout & refused by the reaction (exit status 3)\\
\code{doing\_semi}, kaons, \code{doing\_decay} = 0 & as for pions, but the weight is multiplied by a survival probability that is never computed (0): every weight is 0 & refused by the reaction (exit status 3)\\
\end{simctable}

\begin{simcfindingbox}
\textbf{The phase-space switch outside hydrogen.} From the code (not run): (1) on deuterium the residual mass is
$M_p+M_n-M_N$ with $M_N=0$ \fsrc{dbase.f:447-451}, so the proton is solved against a recoil of mass $M_p+M_n$
\fsrc{event.f:567-579}; (2) for $A\ge3$ every try fails the vertex missing-energy test \fsrc{event.f:403-407, 945}; (3)
with a meson flag the electron energy is never generated \fsrc{event.f:300-316} and keeps the value of the event record
(0 at start-up); (4) with SIDIS kaons and \code{doing\_decay} = 0, the survival probability multiplies the weight
\fsrc{event.f:1563-1564} but is computed only by the SIDIS model, which the phase-space branch does not call
\fsrc{event.f:1449-1451}: the weight is 0. SIMC++ refuses all of these (\tabr{tab:replaces}); the list of Fortran
differences records them as deliberate deviations (the deuterium/heavy and SIDIS refusals) and as a combination that is
not supported yet (the meson flags), see \tabr{tab:findings}.
\end{simcfindingbox}

\subsection{Event flow}
\begin{simctable}{footnotesize}{L{0.05\linewidth}L{0.33\linewidth}L{0.62\linewidth}}
\caption{One try of a phase-space run (one step per row; the general chain is \ctabr{tab:chain}).}\lbl{tab:flow}\\
\toprule\rowcolor{simchead}\simcth{\#} & \simcth{C++ call} & \simcth{what is specific to this run}\\\midrule\endfirsthead
\toprule\rowcolor{simchead}\simcth{\#} & \simcth{C++ call} & \simcth{what is specific to this run}\\\midrule\endhead
\bottomrule\endlastfoot
1 & \code{StandardBeam::generate} & nothing\\
2 & \code{HydrogenElastic::generate} & the two electron slopes only (\eqr{eq:genslopes})\\
3 & \code{EepReactionBase::completeEv} & elastic $E'$ and proton (\eqsr{eq:eprime}{eq:proton}); the phase-space branch is not reached\\
4 & \code{SimcRadiator::radiate} & radiation stays on; photon limit with $M_N=0$ (\eqr{eq:egammatot})\\
5 & \code{MonteCarlo::run} & nothing\\
6 & \code{EepReactionBase::completeRecon} & the reconstructed proton momentum is replaced by $P_h^0$ (\eqr{eq:recon})\\
7 & \code{EepReactionBase::completeMain}, \code{assembleWeight} & $S=1$, $\sigma=1$ (\eqr{eq:weight})\\
8 & \code{Normalise} & normfac $=1$ (\eqr{eq:normfac})\\
\end{simctable}

% =====================================================================================================================
\section{Kinematics, generated variables and generation volume}\lbl{sec:gen}
Only the electron slopes are drawn, flat in the electron edges:
\begin{feq}{reactions/HydrogenElastic.cpp: HydrogenElastic::generate}{event.f:276-277}
y'_e = y'_{e,\mathrm{min}} + \xi_1\,(y'_{e,\mathrm{max}} - y'_{e,\mathrm{min}}),\qquad
x'_e = x'_{e,\mathrm{min}} + \xi_2\,(x'_{e,\mathrm{max}} - x'_{e,\mathrm{min}}).
\lbl{eq:genslopes}
\end{feq}
The limits are those of H(e,e\pr p) except the photon limit (\eqr{eq:egammatot}). The generation volume is the electron
solid angle; it is printed in the \texttt{.hist} (\code{genvol}) but does not enter the normalisation:
\begin{feq}{reactions/EepReactionBase.cpp: EepReactionBase::genVolume}{simc.f:383, 390}
V_\mathrm{gen} = \Delta\Omega_e = (y'_{e,\mathrm{max}}-y'_{e,\mathrm{min}})(x'_{e,\mathrm{max}}-x'_{e,\mathrm{min}})\quad[\mathrm{sr}].
\lbl{eq:genvol}
\end{feq}

% =====================================================================================================================
\section{Completing the event}\lbl{sec:complete}
\subsection{Vertex}
The hydrogen-elastic branch comes first in \code{completeEv}, so the electron energy and the proton are elastic:
\begin{feqa}{reactions/EepReactionBase.cpp: EepReactionBase::completeEv}{event.f:518-522}
E' &= \frac{E_\mathrm{in}\,M_h}{M_h + E_\mathrm{in}(1-\cos\theta_e)}\quad(M_h=M_p),\qquad \text{fail if }E'>E_\mathrm{in},
\lbl{eq:eprime}
\end{feqa}
\begin{feq}{reactions/EepReactionBase.cpp: EepReactionBase::completeEv}{event.f:546-562}
\hat p = \hat q,\qquad P_h = |\vec q|,\qquad E_h = \sqrt{P_h^2+M_h^2},\qquad E_m = P_m = 0.
\lbl{eq:proton}
\end{feq}
The elastic relation uses the hadron mass $M_h$, which is $M_p$ here. The dedicated phase-space branch of
\code{completeEv}, reached only when no other branch applies (a meson flag cleared by the switch, \tabr{tab:replaces}),
puts the hadron at the central momentum:
\begin{feq}{reactions/EepReactionBase.cpp: EepReactionBase::completeEv}{event.f:709-716}
P_h = P_h^0,\qquad E_h = \sqrt{M_h^2+P_h^2},\qquad \delta_h = 0.
\lbl{eq:phspbranch}
\end{feq}

\subsection{Reconstruction}
After the spectrometers the reconstructed hadron momentum is overwritten by the central one, for every phase-space
event:
\begin{feqa}{reactions/EepReactionBase.cpp: EepReactionBase::completeRecon}{event.f:1134-1138, 1347-1349}
P_h^\mathrm{rec} &= P_h^0,\qquad E_h^\mathrm{rec} = \sqrt{M_h^2+(P_h^0)^2},\nonumber\\
E_m^\mathrm{rec} &= \nu^\mathrm{rec}+M_p-E_h^\mathrm{rec},\qquad
\vec P_m^\mathrm{rec} = P_h^0\,\hat p^\mathrm{rec}-|\vec q^\mathrm{rec}|\,\hat q^\mathrm{rec},
\lbl{eq:recon}
\end{feqa}
so the tree columns \code{Em}, \code{Pm} and the $\vec P_m$ components of a phase-space run use $P_h^0$, not the
reconstructed proton momentum; the reconstructed $\delta_h$ (column of the hadron arm) is not changed.

% =====================================================================================================================
\section{Jacobians}\lbl{sec:jac}
As for H(e,e\pr p), the only factor is that of the electron slopes (\ceqr{eq:jacomega}):
\begin{feq}{reactions/EepReactionBase.cpp: EepReactionBase::completeEv}{event.f:498, 1015-1016}
J = \bigl(1+x_e'^2+y_e'^2\bigr)^{-3/2}.
\lbl{eq:jac}
\end{feq}

% =====================================================================================================================
\section{Cross-section model}\lbl{sec:xsec}
\subsection{What the run requires from a model}
Nothing: the run's physics model (\code{sigep} by default) is created but never called, and a JSON choice of another
model has no effect. The spectral-function weight is 1 and the cross section is 1:
\begin{feq}{reactions/EventChain.cpp: assembleWeight}{event.f:1400-1401, 1449-1451, 1561-1562}
S = 1,\qquad \sigma = 1,\qquad w = J\cdot w_\mathrm{gen},
\lbl{eq:weight}
\end{feq}
(no Coulomb factor on hydrogen; on a heavier target the factor $(1+V_C/E_b)^2$ of \ceqr{eq:coulombfocus} would multiply
$\sigma=1$). The tree columns \code{sigcc} = 1 and \code{Weight} = $w$; the \texttt{.hist} line \code{CENTRAL.sigcc} is 1.

\subsection{Normalisation}
The normalisation factor of \ceqr{eq:normfac} is replaced by 1:
\begin{feq}{driver/RunReport.cpp: Normalise}{simc.f:404-405, 930}
\mathrm{normfac} = 1,\qquad \mathrm{wtcontr} = \frac{1}{N_\mathrm{event}}\sum_{\text{contributing}}J_i\,w_{\mathrm{gen},i},
\lbl{eq:normfac}
\end{feq}
so \texttt{wtcontr} is the mean weight of the contributing events per success, without luminosity, generation volume or
number of tries. An acceptance follows from the outputs: since each try draws the electron slopes uniformly in
$V_\mathrm{gen}$ and $J$ converts slope area to solid angle, the electron solid angle accepted by both arms and the cuts is
$\Delta\Omega_\mathrm{acc}=V_\mathrm{gen}\sum_\text{contributing}J_i/N_\mathrm{tried}$ without radiation (with radiation each
event also carries its radiative weight). This interpretation is not printed by the code.

% =====================================================================================================================
\section{Radiation, energy loss and decays}\lbl{sec:rad}
Only the radiation input \code{rad\_flag} is set to 0 (the ``best'' formulas, \csecr{sec:radiation}); with
\code{using\_rad} = 1 the three tails are radiated as for H(e,e\pr p), with the tail-1 limits of note 01 (which use
$M_p$ explicitly). The largest photon energy, however, is computed with the struck-nucleon and recoil masses that the
switch left at 0:
\begin{feq}{setup/SetupLimits.cpp: BuildLimits}{init.f:411-425}
E_\gamma^\mathrm{tot,max} = E_b^\mathrm{max} + M_N - M_{\mathrm{rec},N} - E'_\mathrm{min} - E_{h,\mathrm{min}}
= E_b^\mathrm{max} - E'_\mathrm{min} - E_{h,\mathrm{min}}\quad(M_N=M_{\mathrm{rec},N}=0),
\lbl{eq:egammatot}
\end{feq}
one proton mass below the H(e,e\pr p) value. At high beam energy it becomes negative and no tail can radiate, so every
try fails (the Fortran reference run at 6.4 GeV has 0 events, as SIMC++). Set \code{using\_rad} = 0 for a pure acceptance
run. Energy loss and multiple scattering are those of the proton (\csecr{sec:eloss}); nothing decays.

% =====================================================================================================================
\section{Inputs and outputs}\lbl{sec:io}
\subsection{Inputs}
\begin{simctable}{footnotesize}{L{0.24\linewidth}L{0.12\linewidth}L{0.08\linewidth}L{0.56\linewidth}}
\caption{Inputs specific to a phase-space run (one input per row; the kinematic inputs are those of note 01).}\lbl{tab:inputs}\\
\toprule\rowcolor{simchead}\simcth{input} & \simcth{unit} & \simcth{default} & \simcth{meaning here}\\\midrule\endfirsthead
\toprule\rowcolor{simchead}\simcth{input} & \simcth{unit} & \simcth{default} & \simcth{meaning here}\\\midrule\endhead
\bottomrule\endlastfoot
\code{doing\_phsp} & flag & 0 & 1: $\sigma=1$, normfac $=1$\\
\code{targ\%A} & -- & 0 & must be 1 (hydrogen)\\
\code{doing\_pion} & flag & 0 & must be 0\\
\code{doing\_kaon} & flag & 0 & must be 0\\
\code{doing\_delta} & flag & 0 & must be 0\\
\code{doing\_semi} & flag & 0 & must be 0\\
\code{doing\_rho} & flag & 0 & must be 0\\
\code{doing\_nuc\_elast} & flag & 0 & must be 0\\
\code{using\_rad} & flag & 0 & radiation stays on when 1\\
\code{rad\_flag} & -- & 0 & forced to 0\\
\code{spect\_mode} & -- & 0 & which arms are simulated (e.g.\ the acceptance of one arm only; user manual, spectrometers)\\
\end{simctable}

\subsection{Outputs}
The tree has the (e,e\pr p)-type layout of note 01 (53 columns, 46 in the legacy layout), with \code{sigcc} = 1,
\code{Weight} = $J\,w_\mathrm{gen}$, and \code{Em}, \code{Pm}, \code{Pmx}, \code{Pmy}, \code{Pmz}, \code{PmPar},
\code{PmPer}, \code{PmOop} computed with $P_h^0$ (\eqr{eq:recon}). The \texttt{.hist} reports normfac and \texttt{wtcontr}
of \eqr{eq:normfac}.

% =====================================================================================================================
\section{Validation in SIMC++}\lbl{sec:validation}
\begin{simctable}{footnotesize}{L{0.44\linewidth}L{0.21\linewidth}L{0.35\linewidth}}
\caption{CTest tests that cover phase-space runs (one test per row).}\lbl{tab:tests}\\
\toprule\rowcolor{simchead}\simcth{test} & \simcth{labels} & \simcth{what it checks}\\\midrule\endfirsthead
\toprule\rowcolor{simchead}\simcth{test} & \simcth{labels} & \simcth{what it checks}\\\midrule\endhead
\bottomrule\endlastfoot
\code{setup\_hist\_acc\_phsp\_A} & unit & the set-up against the Fortran \texttt{.hist} (limits, masses)\\
\code{hist2json\_acc\_phsp\_hms\_shms\_A} & unit & the \texttt{.hist} reader on the Fortran run A\\
\code{hist2json\_acc\_phsp\_hms\_shms\_B} & unit & the same for run B\\
\code{l1\_event\_capture\_phsp} & l1, slow & every routine of captured events\\
\code{l1\_event\_capture\_phsp\_tries} & l1, slow & the same for a run with a number of tries\\
\code{l1\_montecarlo\_replay\_phsp} & l1, replay & the spectrometer chain of the capture\\
\code{l1\_montecarlo\_replay\_phsp\_tries} & l1, replay & the same for tries\\
\code{l1\_simcrun\_replay\_phsp} & l1, replay & the complete run of the capture, bitwise\\
\code{l1\_simcrun\_replay\_phsp\_tries} & l1, replay & the same for tries\\
\code{cli\_json\_vs\_inp\_phsp} & cli, l1 & the same run from a \texttt{.inp} and from JSON\\
\code{regression\_identity\_phsp} & regression & compat-mode outputs against the stored baseline\\
\code{regression\_corrected\_phsp} & regression & default-mode outputs against the stored baseline\\
\code{regression\_shipped\_p15f1\_phsp\_sm1\_A} & regression\_shipped, slow & phase space with one arm simulated\\
\end{simctable}
Recorded in the validation ledger (\code{project\_control/VALIDATION\_LEDGER.md}): the statistical comparison of SIMC++
with the Fortran run \code{acc\_phsp\_hms\_shms\_A} fails 0 of 209 comparisons (2026-10-01); the run at 6.4 GeV
(\code{acc\_phsp\_K2\_A}) has 0 events in both programs; the one-arm run \code{p15f1\_phsp\_sm1\_A} differs in the default
mode only by \texttt{wtcontr} ($-24.4$\,\%, the hadron $\delta$ cut of \code{pass\_cuts}). The probes of
\tabr{tab:replaces} were made for this note with the SIMC++ build of commit 8260761 (the code is unchanged at cab5a4d),
the input of \code{acc\_phsp\_hms\_shms\_A} with $-20000$ tries, seed 4242: hydrogen 978 contributing events,
\texttt{wtcontr} $=0.0370$; the other rows as stated.

% =====================================================================================================================
\section{Fortran origin and Fortran findings}\lbl{sec:origin}
\begin{simctable}{footnotesize}{L{0.42\linewidth}L{0.58\linewidth}}
\caption{Fortran routine $\to$ C++ function (one routine per row; the H(e,e\pr p) routines are in note 01).}\lbl{tab:map}\\
\toprule\rowcolor{simchead}\simcth{Fortran (upstream 59ccdf6)} & \simcth{C++ (src/...)}\\\midrule\endfirsthead
\toprule\rowcolor{simchead}\simcth{Fortran (upstream 59ccdf6)} & \simcth{C++ (src/...)}\\\midrule\endhead
\bottomrule\endlastfoot
\file{dbase.f:231}-238 (flags) & \csrc{core/RunSetup.cpp: DeriveReactionFlags}\\
\file{dbase.f:231}-232 (\code{rad\_flag} = 0) & \csrc{setup/SetupReaction.cpp: BuildReactionAndKinematics}\\
\file{dbase.f:887}-888, 952 (messages) & \csrc{setup/SetupBuilder.cpp}\\
\file{event.f} generate, l.217-414 & \csrc{reactions/HydrogenElastic.cpp: HydrogenElastic::generate}\\
\file{event.f} complete\_ev, l.709-716 & \csrc{reactions/EepReactionBase.cpp: EepReactionBase::completeEv}\\
\file{event.f} complete\_recon\_ev, l.1134-1138 & \csrc{reactions/EepReactionBase.cpp: EepReactionBase::completeRecon}\\
\file{event.f} complete\_main, l.1400, 1449-1451 & \csrc{reactions/EventChain.cpp: assembleWeight}\\
\file{init.f:411}-425 & \csrc{setup/SetupLimits.cpp: BuildLimits}\\
\file{simc.f:404} & \csrc{driver/RunReport.cpp: Normalise}\\
\file{simc.f:795}-796 (\texttt{.hist} title) & \csrc{core/RunSetup.cpp: ReactionTitle}\\
\end{simctable}

\begin{simcfindings}{footnotesize}{L{0.2\linewidth}L{0.5\linewidth}L{0.3\linewidth}}
\caption{Fortran findings that touch phase-space runs (one finding per row).}\lbl{tab:findings}\\
\toprule\rowcolor{simchead}\simcth{id} & \simcth{Fortran behaviour} & \simcth{default / compat}\\\midrule\endfirsthead
\toprule\rowcolor{simchead}\simcth{id} & \simcth{Fortran behaviour} & \simcth{default / compat}\\\midrule\endhead
\bottomrule\endlastfoot
\finding{D-009 elastic in phsp} & on hydrogen the elastic kinematics are kept & kept (open question)\\
\finding{D-009 radiation in phsp} & \code{using\_rad} is not switched off & kept (open question)\\
\finding{D-028 F1} & $M_N=0$: the photon limit is one proton mass too low; 0 events at 6.4 GeV & kept (open question)\\
\finding{D-008 phsp} & reconstructed hadron momentum $=P_h^0$, $\sigma=1$, normfac $=1$ & kept (design)\\
\finding{D-006 pass\_cuts} & hadron $\delta$ cut with the electron upper limit & hadron limit / Fortran\\
\finding{D-028 F2} & a NaN proton $y'$ by rounding & $y'=0$ / NaN\\
\finding{U-DEV1} & \code{doing\_nuc\_elast} combined with \code{doing\_phsp} & refused in both\\
\finding{PHYS.2C phsp D/A} & deuterium and heavier targets run with $M_N=0$ & refused in both (exit status 2)\\
\finding{PHYS.2C phsp SIDIS} & SIDIS generation with $\sigma=1$ (kaon weights 0) & refused in both (exit status 3)\\
\finding{PHYS.2C phsp mesons} & meson flags cleared, target class kept & not supported: no layout / no event\\
\end{simcfindings}
\end{document}
